\chapter{Conclusion}\label{ch:conclusion}

The project set out to let a user ask questions over a collection of scientific papers and receive answers whose
every claim can be traced to a passage and a page, with an explicit refusal when the papers do not support an
answer. The resulting system ingests papers from arXiv and from uploaded PDFs, preserves provenance from page to
chunk, retrieves passages by combining dense and lexical search with cross-encoder re-ranking, and generates answers
with a local model whose citations are resolved on the server.

The principal engineering decisions were to keep all inference local and refuse any hosted fallback, to let the
model cite only passage labels so that the server can verify every citation, to persist evidence snapshots so that
history survives changes to the corpus, and to treat evaluation as part of the system: a versioned evaluation set,
reproducible run records, and a regression gate in continuous integration.

The verification evidence is mixed in a way the report has tried to state precisely. Citation integrity and
provenance were verified deterministically across every stored answer and every stored chunk, and the default
retrieval configuration reached a Recall@5 of 0.875 on the evaluation set. Against that, the latency target was not
met, one in ten unanswerable questions received a wrong-paper answer, spot checks found claims not supported by their
cited passage, and the evaluation labels have not yet been reviewed by a person.

For the course, the project exercised the full software-engineering cycle: requirements with stable identifiers,
architecture decisions recorded with their alternatives, design modelled in UML, layered automated testing, security
and reproducibility reviews, and an acceptance decision based on evidence rather than a successful build. For the
institute, it provides a documented, locally deployable base that future students can extend and evaluate with the
same tools.
